\section{潜在应用}\label{sec:app}
随着网络作为数据建模结构在社会、自然和学术领域中表示多样化关系信息的广泛应用，图提示已开始从通用自适应技术走向具体的应用场景。在本节中，我们主要关注在图结构数据或图表示上实例化的提示，同时也会注意到图知识支持基于提示推理的相关设置。我们讨论其在在线社交网络、推荐系统、知识管理和生物学中的新兴应用，重点关注动态交互图、用户-物品图、知识图谱和生物医学图。


\textbf{在线社交网络。}
在线社交平台由用户组成，用户可以表示为节点，他们的社交连接、转帖痕迹或互动事件形成在线社交网络（OSNs）。先前研究已调查了提示机制在识别OSNs中虚假信息以防止恶意攻击方面的潜力。例如，\citet{wu2023promptandalign} 将文本提示应用于预训练语言模型（PLMs）以提炼通用语义信息，然后将其与社交网络中的传播动态相结合以改进分类。除了面向文本的提示，OSNs也可以被视为节点状态和边随时间演变的动态交互图。\citet{chen2024prompt} 表明提示学习可用于将时间交互表示适应下游任务，DyGPrompt \cite{yu2025dygprompt} 进一步设计节点-时间条件提示以桥接动态图预训练和下游推理之间的任务-目标差距和时间-分布差距。这些进展使OSN应用不那么投机取巧：少样本虚假信息检测、恶意账户识别和交互级风险评估都可以从提示中受益，这些提示在标记数据稀缺且获取成本高昂时编码时间和结构上下文。


\textbf{推荐系统。}
电子商务平台为利用推荐系统增强在线服务提供了宝贵机会，因为用户-物品交互自然形成二分图或异构图。早期图提示调优研究探索了跨域推荐和对比预训练自适应 \cite{wu2023personalized, hao2024motifbased}。对于跨域推荐，\citet{yi2023contrastive} 在将预训练推荐模型迁移到目标域时引入额外提示节点，从而用有限的目标监督改进域自适应。同时，\citet{yang2023empirical} 提出个性化用户提示以弥合对比代理任务 \cite{wu2021selfsupervised, yu2022are} 和下游推荐目标之间的差距。近期工作进一步将这一方向从静态自适应扩展到动态和流式用户-物品图。GraphPro \cite{yang2024graphpro} 结合图预训练与时间和结构提示以改进动态推荐，而 GPT4Rec \cite{zhang2024gpt4rec} 引入节点级、结构级和视图级提示以指导演化交互图上的流式推荐。这些研究表明，图提示特别适合推荐系统，在这些系统中用户兴趣、物品热门度和交互模式持续变化，并且需要为冷启动、跨域和持续推荐设置进行高效自适应。


\textbf{知识管理。} 知识图谱（KGs）提供了另一个可以直接应用提示调优到图结构知识的应用场景。一条核心工作是使用预训练KG模型对KGs进行提示以促进不同KG数据和任务之间的知识迁移。在这条工作中，\citet{zhang2023structure} 设计了结构感知预训练目标，然后采用提示调优将任务相关知识迁移到下游KG任务。近期KG补全研究进一步表明，自适应或双重提示可以帮助将结构化关系模式与任务特定预测目标对齐：APKGC \cite{maoliniyazi2026apkgc} 使用自适应提示进行KG补全，FFCL-KGC \cite{zhu2026ffclkgc} 将双重提示纳入特征融合对比学习。同时，一个相关但更广泛的方向将来自KGs的基础知识与LLMs的推理和语言生成能力相结合 \cite{wang2023knowledge, robinson2023leveraging, park2023graphguided, zhang2023benchmarking, pan2023unifying}。例如，\citet{tian2023graph} 引入图神经提示以时间和参数高效的方式将KG知识适应LLMs。这些研究表明，图提示可以支持KGs内的结构化知识迁移，而KG基础的LLM提示为图增强知识推理提供了相邻方向。


\textbf{生物学。} 分子可以自然地表示为图，其中原子作为节点，化学键作为边 \cite{rong2020selfsupervised}。这种图建模为将图表示学习方法应用于分子性质预测、药物相互作用分析和蛋白质建模等任务提供了基础 \cite{rong2020selfsupervised, liu2023molca, zhao2023gimlet, edwards2023synergpt}。早期分子图学习方法的训练通常针对单个数据集的任务特定模型 \cite{rong2020selfsupervised, sun2020infograph}，这限制了跨任务泛化。虽然LLMs和语言模型已被探索为理解分子的通用工具 \cite{qian2023can, liu2023one}，但纯文本描述可能忽略分子的内部图结构 \cite{zhang2023large}。近期图-语言模型试图在少样本或零样本设置下保留结构依赖和语义知识 \cite{zhao2023gimlet, edwards2023synergpt}。例如，GIT-Mol \cite{liu2023gitmol} 使用LoRA \cite{hu2021lora} 在分子LLM系统中进行高效下游自适应。图提示现在提供了更直接的生物医学证据。MOAT \cite{long2024moat} 为3D分子图设计原子级、几何级和任务级提示；DDIPrompt \cite{wang2024ddiprompt} 将图提示学习应用于稀有和不平衡的药物-药物相互作用事件预测；Prompt-DDG \cite{wu2024promptddg} 使用微环境感知层次提示来预测蛋白质-蛋白质相互作用中的突变效应。与蛋白质多聚体提示学习 \cite{gao2024protein} 结合，这些研究表明图提示可以在分子、药物安全和蛋白质相互作用应用中支持跨任务的高效自适应。
